% !TeX root = ../sections/02_exercises.tex

\begin{exercise}
	\(\mathbb{Z}\) のイデアル \(3\mathbb{Z}\) と \(5\mathbb{Z}\) について，
	中国式剰余定理から
	\[
	\mathbb{Z}/15\mathbb{Z}
	\simeq
	\mathbb{Z}/3\mathbb{Z}
	\oplus
	\mathbb{Z}/5\mathbb{Z}
	\]
	が成り立つことを踏まえて，次の問いに答えよ。
	
	\begin{enumerate}
		\setcounter{enumi}{-1}
		
		\item
		\[
		1=3a+5b
		\]
		を満たす整数 \(a,b\) を一組求めよ。
		
		\item
		\(\mathbb{Z}/3\mathbb{Z}\oplus\mathbb{Z}/5\mathbb{Z}\) において，
		\[
		(5b,5b)=(1,0)
		\]
		であることを示せ。
		
		\item
		\(\mathbb{Z}/3\mathbb{Z}\oplus\mathbb{Z}/5\mathbb{Z}\) において，
		\[
		(3a,3a)=(0,1)
		\]
		であることを示せ。
		
		\item
		\(\mathbb{Z}/3\mathbb{Z}\oplus\mathbb{Z}/5\mathbb{Z}\) において，
		\[
		(4\cdot 3a+2\cdot 5b,\,
		4\cdot 3a+2\cdot 5b)
		=
		(2,4)
		\]
		であることを示せ。
		
		\item
		\(3\) で割った余りが \(2\) となり，
		\(5\) で割った余りが \(4\) であるような整数を求めよ。
	\end{enumerate}
\end{exercise}

\begin{background}
	この問題の中心は，中国式剰余定理である。
	互いに素な整数 \(3,5\) に対して，
	\[
	\mathbb{Z}/15\mathbb{Z}
	\simeq
	\mathbb{Z}/3\mathbb{Z}
	\oplus
	\mathbb{Z}/5\mathbb{Z}
	\]
	が成り立つ。
	この同型は，整数 \(n\) に対して
	\[
	n \longmapsto (n \bmod 3,\ n \bmod 5)
	\]
	を対応させる写像によって与えられる。
	
	また，\(3\) と \(5\) は互いに素なので，ベズーの等式より
	整数 \(a,b\) が存在して
	\[
	1=3a+5b
	\]
	と書ける。
	
	この等式から，
	\[
	5b \equiv 1 \pmod{3},
	\qquad
	5b \equiv 0 \pmod{5}
	\]
	および
	\[
	3a \equiv 0 \pmod{3},
	\qquad
	3a \equiv 1 \pmod{5}
	\]
	が従う。
	したがって，\(5b\) は第1成分だけを取り出す元，
	\(3a\) は第2成分だけを取り出す元として働く。
\end{background}

\begin{strategy}
	まず，ベズーの等式
	\[
	1=3a+5b
	\]
	を満たす整数 \(a,b\) を具体的に求める。
	これは互除法，または簡単な代入で見つければよい。
	
	\begin{enumerate}
		\item[(0)]
		\(3a+5b=1\) を満たす整数の組を一つ求める。
		この組は一意ではないが，一組あれば以降の計算に使える。
		
		\item[(1)]
		\(1=3a+5b\) を法 \(3\) で見る。
		すると
		\[
		5b\equiv 1 \pmod{3}
		\]
		が分かる。
		また，法 \(5\) では
		\[
		5b\equiv 0 \pmod{5}
		\]
		である。
		よって \((5b,5b)=(1,0)\) となる。
		
		\item[(2)]
		同様に，\(1=3a+5b\) を法 \(5\) で見る。
		すると
		\[
		3a\equiv 1 \pmod{5}
		\]
		が分かる。
		一方，法 \(3\) では
		\[
		3a\equiv 0 \pmod{3}
		\]
		である。
		よって \((3a,3a)=(0,1)\) となる。
		
		\item[(3)]
		(1)，(2) で得た
		\[
		5b \leftrightarrow (1,0),
		\qquad
		3a \leftrightarrow (0,1)
		\]
		を使う。
		つまり，
		\[
		4\cdot 3a+2\cdot 5b
		\]
		は，第1成分では \(2\cdot 5b\)，
		第2成分では \(4\cdot 3a\) が効くと考える。
		
		\item[(4)]
		余りの条件
		\[
		n\equiv 2 \pmod{3},
		\qquad
		n\equiv 4 \pmod{5}
		\]
		を，中国式剰余定理を使って一つの整数にまとめる。
		(1)，(2) の元を用いて，
		\[
		n=2\cdot 5b+4\cdot 3a
		\]
		とおけば，指定された二つの合同条件を同時に満たす。
	\end{enumerate}
\end{strategy}